\section{SCIP建模与求解}
\label{sec:03_01}

第二章确定了候选任务、资源占用关系与收益，本节说明如何将这些对象交给求解器，以及如何解释返回的选择方案、上下界和停止状态。SCIP是公开的通用优化求解器，其框架结合了预处理、分支定界、割平面和启发式等机制\cite{SCIP100}，并非本书提出的超图算法。历史组合实验使用SCIP 10.0.2、Python接口PySCIPOpt 6.2.1及线性规划求解器SoPlex 8.0.2。本节的求解对象限于第二章的纯组合模型，后续巡检联合模型的运行状态另行讨论。

\subsection{从资源关联到整数规划模型}

沿用第二章的定义，$V$为资源集合，$E$为候选任务索引集合。对每个任务$e\in E$，非空集合$R_e\subseteq V$列出其必须共同占用的资源；选择该任务时，对$R_e$中的每种资源均占用一份。不同任务可以具有相同的资源集合，仍须保留各自的任务索引。记$n=|E|$为候选任务数，$m=|V|$为资源数。任务收益$w_e\in\mathbb{R}_{>0}$为无量纲正数，资源容量$b_v\in\mathbb{Z}_{\geq0}$表示资源$v\in V$允许的最大占用次数。

二元变量$x_e\in\{0,1\}$表示是否完整选择任务$e$。关联参数$A_{ve}\in\{0,1\}$在$v\in R_e$时为$1$，否则为$0$；已有邻接任务集合$N(v)=\{e\in E:v\in R_e\}$列出所有需要资源$v$的任务。交给SCIP的仍是第二章的同一个整数规划，其紧凑形式为
\begin{equation}
 z^{\star}=\max_{x\in\mathcal{F}}\sum_{e\in E}w_ex_e,
 \qquad
 \mathcal{F}=\left\{x\in\{0,1\}^{n}:
 \sum_{e\in N(v)}x_e\leq b_v,
 \ \forall v\in V\right\}.
\label{s0301:submitted_model}
\end{equation}
其中$\mathcal{F}$是允许的任务组合，$z^{\star}$是最大无量纲收益。每个资源行左端为占用该资源的已选任务数，与右端容量具有相同的计数含义。全部当前组合实验均取$b_v=1$，因此同一资源不能被两个已选任务重复占用。

建模时，先按原任务标识创建$n$个二元变量，再按资源标识创建$m$条容量约束，最后指定收益最大化目标。约束表达式仅包含实际关联的变量。记关联总数为
\begin{equation}
 p=\sum_{v\in V}|N(v)|
   =\sum_{e\in E}|R_e|,
 \qquad p\in\mathbb{Z}_{\geq0}.
\label{s0301:incidence_count}
\end{equation}
$p$也是原容量矩阵的非零元数。压力测试程序先遍历任务的资源列表建立$N(v)$，再生成各资源行，避免存储稠密的$m\times n$关联矩阵。早期十二任务脚本逐资源扫描任务列表，生成的约束同样只包含实际关联的变量。该算例对应$12$个二元变量、$12$条原资源容量约束和$42$个非零关联。

原资源行在提交模型中保留，SCIP随后可以在内部预处理中变换模型。在单位容量下，对共享资源的任务对建立冲突约束，也能表达相同的二元可行域；本批SCIP建模直接采用原资源行。因此，$|R_e|$较大说明单个任务需要较多种资源，本身不构成不可约高阶耦合或求解性能优势的证据。

\subsection{可行下界、松弛上界与停止精度}

对最大化问题，一组满足原约束的二元选择$\hat{x}$提供可实现的收益。记当前整数可行解收益为$L$，有效上界为$U$，则
\begin{equation}
 L=\sum_{e\in E}w_e\hat{x}_e
 \leq z^{\star}\leq U,
 \qquad \hat{x}\in\mathcal{F}.
\label{s0301:bounds}
\end{equation}
$L$、$U$和$z^{\star}$均为无量纲收益。下界表示已经找到的组合能够实现多少收益，上界则限制尚未排除的最优收益，两者承担不同的证据功能。

这种关系可以从分支定界过程理解。在某个搜索节点，部分变量的取值范围已由分支决定，将二元条件放松为$0\leq x_e\leq1$，得到包含该节点全部整数可行解的线性松弛。该松弛的有效最优值构成节点整数收益的上界。对取分数值的变量继续分支，会把整数候选分配到更小的搜索区域；找到满足整数性和原约束的解，则可以提高$L$。若某节点不可行，或其有效上界不超过已有$L$，该节点即可排除。全局$U$必须覆盖尚未排除的所有搜索区域，任意单个节点的松弛值未必具有这一性质。预处理、割平面和内部启发式可以推进搜索，但不改变式\eqref{s0301:bounds}中的上下界方向。

由于容量非负，全零选择总是可行解；由于权重为正且每个任务至多选择一次，搜索开始前已有
\begin{equation}
 0\leq z^{\star}\leq\sum_{e\in E}w_e.
\label{s0301:trivial_bounds}
\end{equation}
压力测试的初始轨迹记录的正是空选择及收益总和上界。这些信息无需搜索即可得到，并不表示已经发现了非空方案。

当上下界均为有限值，且$L>0$、$U\geq L$时，相对差为
\begin{equation}
 g=\frac{U-L}{L},\qquad
 g\leq0.05\ \Longleftrightarrow\ U\leq1.05L.
\label{s0301:positive_gap}
\end{equation}
$g$无量纲，其分母是当前可行收益。程序保存求解器返回的差距字段，本批$5\%$目标在上述条件下按式\eqref{s0301:positive_gap}解释。若上下界有效，达到该目标还意味着
\begin{equation}
 \frac{L}{z^{\star}}\geq\frac{L}{U}\geq\frac{1}{1.05}.
\label{s0301:gap_implication}
\end{equation}
这给出当前可行解相对最优收益的保证，但不要求$L$已经等于$z^{\star}$。

零目标须单独处理。若$L=0<U$，可报告绝对差$U-L$及尚未取得正收益可行解的事实，相对差式\eqref{s0301:positive_gap}不适用；若有效界满足$L=U=0$，则最优收益为零，$0/0$仍没有定义。无穷界、缺失界或符号不同的上下界也不适用该正值公式。本模型已有空选择这一可行解，故有效下界至少为零；接口暂时返回的其他数值需要结合哨兵值和求解阶段辨认。

同一有效上界还可以评价独立核验的启发式解。设其收益为$h\geq0$，且$z^{\star}>0$、$U>0$，则
\begin{equation}
 h\leq z^{\star}\leq U
 \quad\Longrightarrow\quad
 \frac{h}{U}\leq\frac{h}{z^{\star}}\leq1.
\label{s0301:heuristic_ratio}
\end{equation}
因此，$h/U$是该解相对于真实最优收益的比值下界。当$L>0$时，$h/L$仅比较启发式收益与当前SCIP可行收益；若启发式找到了更好的方案，该值可以超过$1$。两种比值的含义不同，且均要求参与比较的收益与上界来自同一实例、同一目标和同一可行域。

上述数学关系还需与浮点求解记录的证据强度区分。历史脚本将SCIP返回的\texttt{optimal}状态保存为最优性标志，它表示求解器在当次数值设置下报告最优。所给代码没有提供定向舍入的界证书或有理数精确验证过程，因此该状态属于浮点求解结论。SCIP 10具有精确求解模式，也不意味着本批实验开启了该模式。对原输入独立核验选择方案能够支持可行下界，却不会同时完成上界的独立证明。

\subsection{不同实验批次的启动与停止条件}

求解器名称相同，各批次的初始信息和计算预算仍可能不同。已有记录需要按以下条件解释。

\begin{enumerate}
 \item 十二候选任务算例设置单次求解时间上限为$10\,\mathrm{s}$，最大线程数为$1$，随机种子偏移为$0$。质量运行与$21$次重复测时分开组织。重复测时使用同一个实例，每次重新创建模型，但进程和已加载库保持热状态。这些重复属于同一输入的时间观测，不是独立实例，也不复用上一轮的整数方案。
 \item 初期实验保留了$10\,\mathrm{s}$限制及首次正收益解口径的历史记录。首次取得正收益只说明已经找到非空可行方案。后续以$5\%$差距为目标的独立续跑复用已有可行解作为热启动，其初始信息与初期运行不同；复用选择方案也不等同于恢复先前的完整搜索树。
 \item $59$个压力条件中的SCIP运行均为冷启动，只显式提供全零解，不注入超图方法的选择结果或历史SCIP方案。每次设置最大线程数为$1$、随机种子偏移为$0$、差距目标为$0.05$，不设置求解时间截断；程序检查时间限制参数保持默认值$10^{20}\,\mathrm{s}$。运行仍受内存保护和异常停止影响。
 \item 原作者完整流程所对应的$7$个实例中，前$6$个实例使用历史SCIP解作为对照，\texttt{superblue19}使用独立冷启动SCIP运行。因此，这组对照并非统一重新执行的$7$次冷启动，其时间记录不构成相同预算下的速度排名。
\end{enumerate}

压力批次设有两层内存限制：SCIP内部内存限制为$1536\,\mathrm{MiB}$，外部保护在采样到进程常驻内存超过$2560\,\mathrm{MiB}$时终止进程，其中$1\,\mathrm{MiB}=2^{20}$字节。前者为求解器参数，后者针对整个进程，两者的计数范围不同。外部监控按$1\,\mathrm{s}$间隔采样，因此所记峰值为采样峰值，而非连续监测得到的精确分配峰值。

压力调度允许两个SCIP进程并行，另设一个原生算法执行槽。单次SCIP求解设置为单线程，不等于整台机器同时只有一个求解任务。线程数、进程并发、热启动信息和停止条件共同限定了时间比较的范围。因内存保护、异常退出或缺少完整输出而结束的运行，应保留失败或未完成状态，不能计作达到差距目标。

\subsection{计时边界与重启成本}

为区分模型接口开销和搜索开销，设$t_0$为模型构建开始时刻，$t_1$为调用优化过程前的时刻，$t_2$为优化调用返回时刻，$t_3$为本次输出提取阶段的计时终点。各时刻取自同一单调计时源，差值以秒计。定义
\begin{align}
 T_{\mathrm{build}}&=t_1-t_0,
 &T_{\mathrm{solve}}&=t_2-t_1,
 \label{s0301:timing_build_solve}\\
 T_{\mathrm{output}}&=t_3-t_2,
 &T_{\mathrm{total}}&=t_3-t_0
 =T_{\mathrm{build}}+T_{\mathrm{solve}}+T_{\mathrm{output}}.
 \label{s0301:timing_total}
\end{align}
四项时间均非负。建模时间包含变量、资源行和目标的创建；压力脚本还把空解设置及回调配置计入该阶段。求解时间是优化调用的外部墙钟时间。输出时间覆盖计时终点之前的解提取及已执行的数据整理：十二任务脚本在整理求解状态和界等字段之后取$t_3$，压力脚本则在提取任务标识之后取该时刻。因而这一分段名称并不保证两个脚本具有完全相同的输出操作范围。

原十二任务脚本另存SCIP内部求解时间，它与外部$T_{\mathrm{solve}}$的计时边界不同。诊断文件导出、原JSON上的独立核验及最终结果落盘位于相关计时区间之外，故$T_{\mathrm{total}}$也不是整个实验进程从启动到退出的耗时。压力批次另有进程墙钟时间与CPU时间记录，前者覆盖更宽的进程执行范围，后者计量处理器使用时间，应分别解释。

续跑还需保留产生热启动解的历史成本。若同一实例依次经历$K$次有关尝试，$K\in\mathbb{Z}_{>0}$，令$T_{\mathrm{process}}^{(k)}\geq0$表示第$k$次尝试的进程墙钟时间，$k=1,\ldots,K$，则在这些尝试的记账范围内，累计运行投入为
\begin{equation}
 C_{\mathrm{run}}=\sum_{k=1}^{K}T_{\mathrm{process}}^{(k)}.
\label{s0301:restart_accounting}
\end{equation}
$C_{\mathrm{run}}$以秒计，是各次尝试耗时之和；有并发时，它不等于日历时间跨度，也不等于CPU时间。后一次运行读取已有方案后重新开始计时，先前取得该方案的成本仍然存在。记录若仅覆盖单次续跑，就应标为单次续跑时间。实时回调中的收益、上界和耗时属于求解器报告的中间快照，本节不据尚未核清的快照报告已完成的性能结果。

\subsection{从浮点输出到原输入核验}

优化调用返回后，脚本从当前最好解提取变量值，并以大于$0.5$作为输出任务标识的判据。这一步将数值解转换为待核验的离散选择，资源可行性仍需检查。设输出任务标识通过无重复和存在性检查后形成集合$S\subseteq E$，核验器回到原JSON读取任务收益、资源列表与容量，重新计算
\begin{align}
 q_v(S)&=\sum_{e\in S}A_{ve},
 &&v\in V,
 \label{s0301:verified_load}\\
 W(S)&=\sum_{e\in S}w_e.
 \label{s0301:verified_value}
\end{align}
$q_v(S)\in\mathbb{Z}_{\geq0}$是资源$v$的占用次数，$W(S)\in\mathbb{R}_{\geq0}$是无量纲选择收益。只有$q_v(S)\leq b_v$对每个$v\in V$均成立，输出才满足原资源容量约束。核验器不删除冲突任务来修复方案，以保持被核验对象与算法实际输出一致。

当原始权重为整数时，$W(S)$按原数据进行整数求和，无需浮点可行性容差。资源占用也是整数计数，单位容量约束下，占用次数$2$就是违反容量$1$。求解器浮点目标与$W(S)$是否一致则是另一项检查，其容差仅用于数值报告的一致性，不改变组合可行性的定义。

已有初始解核验记录覆盖$33$个条件、$78$条解记录，其中包含$9$条诊断记录；所检查的选择均满足原输入容量约束。该核验采用$10^{-6}$的绝对容差比较浮点目标与整数收益。例如，\texttt{superblue19}的随机权重条件中，一条较后完成的SCIP记录对应原输入整数收益$11642395$，浮点报告为$11642394.999999993$，浮点报告减去整数收益的差约为$-7.45\times10^{-9}$。该差异位于比较容差内，不应记作收益损失。压力运行脚本的目标一致性断言采用$10^{-5}$的绝对容差，这与初始解核验记录属于不同的核验设置。

上述$78$条记录包含历史结果和较后完成结果，并非$78$个独立重复实验。核验支持输出选择的可行性及原输入收益，不额外证明SCIP上界或最优性。解释一条SCIP结果时，应将停止状态、经原输入核验的任务集合及收益、求解器报告的上下界和差距，与该次运行的启动方式和计时范围对应起来，才能确定其可行性证据与解质量比较的适用范围。